
\chapter{What This Battery-First Thesis Proves and What It Leaves Open}
\label{ch:what-proves}

A dissertation that claims a general principle but validates it in one domain
must be explicit about the boundary between proven and projected.  This
chapter draws that boundary.

\section{What Is Fully Validated in Battery}

The following results are established by reproducible experiment on the
locked CPSBench harness (seed~42, 168-hour episodes, four fault
scenarios, two regions) and backed by formal proofs:

\begin{enumerate}
  \item \textbf{OASG existence.}  The Observation-Action Safety Gap is
    measurably non-zero for all three baseline controllers (deterministic LP,
    CQR-only, robust MPC) under dropout rates as low as~5\%.
    Chapter~\ref{ch:main-results} reports TSVR of 3.9--4.2\% for the
    deterministic LP at 20\% dropout.
  \item \textbf{DC3S effectiveness.}  The DC3S shield achieves exactly
    0\% TSVR across all tested dropout rates $d \in [0, 0.30]$, both
    regions, and all fault combinations.
  \item \textbf{Zero true-state violation rate.}  No single step in any
    experiment produces a true-state SOC breach under DC3S protection.
  \item \textbf{Formal safety bound.}  Theorem~\ref{thm:safety-preservation}
    guarantees that, under exchangeability of the calibration residuals and
    bounded observation error, the true-state violation probability is at
    most $\alpha$ with confidence $1 - \delta$.  The empirical result is
    strictly stronger: violation rate is zero, not merely bounded.
  \item \textbf{Bounded intervention cost.}  The shield's intervention
    rate remains below 3.1\% under \texttt{drift\_combo} and below
    20\% even at extreme 30\% dropout.
\end{enumerate}

\section{What Is Partially Proved}

Several claims are supported by evidence but carry caveats:

\begin{enumerate}
  \item \textbf{Conditional coverage.}  CQR group coverage
    (Chapter~\ref{ch:conditional-coverage}) is validated on the four
    tested fault families.  Coverage under \emph{out-of-distribution}
    faults---sensor types or degradation modes not in the calibration
    set---has not been empirically verified.
  \item \textbf{Aging effects.}  The aging-aware recalibration
    (Chapter~\ref{ch:temporal-extensions}) is demonstrated on synthetic
    degradation curves.  Field data spanning multi-year battery aging
    is not available in the current dataset.
  \item \textbf{Certificate half-life.}  The blackout study
    (Chapter~\ref{ch:cert-half-life-blackout}) now establishes a bounded
    safe-hold result and a conservative width-expansion proxy rather than a
    measured finite safety-breakdown horizon.  Longer horizons (weeks,
    months) remain extrapolated, not measured.
  \item \textbf{Transfer across regions.}  The US--DE transfer stress
    test (Chapter~\ref{ch:transfer-generalisation}) shows positive
    results for two grid regions.  Generalisation to grids with
    fundamentally different market structures (e.g., islanded
    microgrids) is untested.
\end{enumerate}

\section{What Is Framework-Level but Not Yet Validated Universally}

The ORIUS architecture makes structural claims that go beyond the battery
domain.  These claims are implemented and auditable in the repo, but the
non-battery evidence remains shallower than the locked battery proof surface:

\begin{enumerate}
  \item \textbf{Multi-domain portability.}  Chapter~\ref{ch:battery-to-orius}
    argues that the adapter pattern generalises.  Non-battery adapters now
    exist in ORIUS-Bench (navigation, industrial, healthcare, aerospace,
    vehicle), but they should be read as portability and auditability
    evidence rather than as battery-grade proof surfaces.
  \item \textbf{Multi-asset composition.}  The fleet-level shield
    (Chapter~\ref{ch:compositional-safety-fleets}) is validated for a
    two-battery fleet.  Composition across \emph{heterogeneous} asset
    types (battery + solar inverter + load controller) is an open problem.
  \item \textbf{Runtime portability.}  CertOS
    (Chapter~\ref{ch:certos-runtime}) now has bounded systems evidence,
    including lifecycle artifacts, invariants, recovery behavior, and a toy
    second-domain demonstration.  That is runtime evidence, not a universal
    runtime-safety proof.
  \item \textbf{Adversarial completeness.}  The active-probing audit
    summarized in Chapter~\ref{ch:outside-current-evidence} detects the
    tested spoofing attack.  Completeness against all possible adversarial
    strategies is not claimed and likely unachievable in general.
\end{enumerate}

\section{Why Battery-First Is the Correct Sequence}

Starting from the battery domain is not a convenience---it is a research
strategy.  Batteries offer hard state constraints (SOC bounds), high
action-to-state sensitivity, and network-mediated telemetry: exactly the
combination that makes the OASG both dangerous and measurable.  Proving
the framework here means the hardest measurement problem is solved first.
Domains with softer constraints or lower sensitivity are structurally
easier and can be addressed with the same machinery.

\section{What Second Domain Should Come Next}

The most informative second domain is one that \emph{differs} from battery
along the axes identified in Chapter~\ref{ch:battery-to-orius}: nonlinear
dynamics, multi-dimensional state, and non-box constraints.  Building HVAC
thermal management or autonomous vehicle lane-keeping would test
whether the conformal region and projection oracle generalise beyond
scalar-state, linear-dynamics plants.

\section{How to Expand Without Rewriting the Core}

The adapter architecture ensures that the ORIUS core library---OQE scoring,
RAC inflation, certificate issuance, audit trail---is domain-independent
code.  Adding a new domain means writing a new \texttt{DomainAdapter}
implementation and a new CPSBench plant model, not forking the safety
pipeline.  The core's test suite remains valid across all adapters,
providing regression guarantees as the framework grows.

\section{Chapter Summary}

This thesis fully validates the ORIUS principle for the battery dispatch
domain: the OASG exists, DC3S closes it, and the formal safety bound
holds empirically.  It partially validates conditional coverage, aging
recalibration, and regional transfer.  It argues but does not yet prove
that the architecture generalises to other CPS domains.  Being explicit
about this boundary is not a weakness of the battery-first strategy;
it is the point.
